% !PRINTED 101
\weberpage{101}
%<<CONT
verschieden und können nicht durch \gesperrt{Tschirnhausen-Transformation} in einander übergeführt werden, weil sie zu primären Theilern von verschiedenen vollen Kreistheilungskörpern gehören, oder, wenn sie in demselben vollen Kreistheilungskörper enthalten sind, verschiedene Gruppen haben.

Man kann freilich noch aus anderen Einheitswurzeln Kreistheilungsperioden bilden, die zu cubischen Gleichungen führen. Diese Perioden können aber durch niedrigere Einheitswurzeln ausgedrückt werden, und sind nicht primär. So erhält man z.\,B., wenn $r$ eine 35\textsuperscript{ste} Einheitswurzel ist, eine Periode von 8 Gliedern, deren Exponenten aus den Potenzen von $-8$ gebildet sind:

\[
\eta = r + r^{-1} + r^6 + r^{-6} + r^8 + r^{-8} + r^{13} + r^{-13},
\]
die also auch einer cubischen Gleichung genügt.

Es ist aber
\[
1 + r^7 + r^{-7} + r^{14} + r^{-14} = 0,
\]
und wenn man diese Gleichung mit $(r^{15} + r^{-15})$ multiplicirt:
\[
r^{15} + r^{-15} + r + r^{-1} + r^6 + r^{-6} + r^8 + r^{-8} + r^{13} + r^{-13} = 0,
\]
also
\[
\eta = -(r^{15} + r^{-15}),
\]
was unter den Perioden der 7\textsuperscript{ten} Einheitswurzeln schon vorkommt. Diese Erscheinung erklärt sich daraus, dass $m = 35$ nicht die in \S.\,21 verlangte Eigenschaft hat, dass $q - 1$ für alle in $m$ aufgehenden Primzahlen $q$ durch 3 theilbar ist.

\begin{center}
\S.\,24.

\gesperrt{Biquadratische Kreistheilungskörper.}
\end{center}

Bei den biquadratischen Kreistheilungsgleichungen hat man zwei Arten zu unterscheiden. Wir können zwei Invarianten $i_1 = i_2 = 2$ oder nur eine Invariante $i_1 = 4$ annehmen. Wir wollen den ersten Fall voranschicken.

Für $m$ haben wir in diesem Falle nach der Zusammenstellung \S.\,21 folgende Bedingungen:

$m$ kann ungerade sein, durch 4 oder durch 8, aber durch keine höhere Potenz von 2 theilbar sein. Keine ungerade Primzahl kann mehr als einmal in $m$ aufgehen, und es müssen mindestens zwei Indexmoduln vorhanden sein; also müssen ausser in
%CONT>>
